% Original research fragment prepared during the 10 November 2026 iteration.
% Requires amsmath, amssymb, amsthm.  No standalone document preamble.
\section{A scalar successor-seeded recursor is sufficient}
\label{sec:scalar-successor-bootstrap}\label{sec:hybrid-bootstrap}

This section gives an internal predecessor construction for the scalar
successor-seeded recursor.  Its update has only the previous state as a
local variable.  It does not receive the old counter.  The construction
first exchanges the numerical states zero and one to manufacture
zero-start iteration.  It then uses affine, rather than ordinary,
addition in a factor-map construction.  This avoids assuming either a
predecessor or an unrestricted selector during the bootstrap.

\subsection{Primitive syntax and available reasoning}

The numerical terms are variables, $0$, and
\[
 R(n,a;u.h).
\]
Only $u$ in the update $h$ is bound.  All its other variables are passive
parameters.  Every body is fixed finite earlier syntax.  Write
\[
 Sa:=R(0,a;u.0),\qquad 1:=S0,
\]
and postulate precisely the two computation schemes
\begin{align}
 R(0,a;u.h)&=Sa,\tag{R0}\label{ss:r0}\\
 R(Sn,a;u.h)&=h[R(n,a;u.h)/u].\tag{RS}\label{ss:rs}
\end{align}
Thus successor is an abbreviation, and the zero equation for its chosen
defining body is tautological.  No numerical function variable or
uniform interpreter is introduced.

The bootstrap uses equality transport and the following two derived
rules, only for literal equations $E$:
\[
 \frac{\Gamma\vdash E(0)\qquad\Gamma\vdash E(Sz)}
      {\Gamma\vdash E(t)},\qquad
 \frac{\Gamma\vdash E(0)\qquad\Gamma,E(z)\vdash E(Sz)}
      {\Gamma\vdash E(t)}.
 \tag{NC, EI}\label{ss:nc-ei}
\]
Here $z$ is fresh for the retained context and fixed parameters.
These rules follow directly from native DI-case and DI-progress,
respectively, with an unused witness variable, followed by DE and QE.
They therefore precede the reconstruction of ordinary logical
connectives.  Reflexivity is also available before any selector:
\eqref{ss:rs} gives $R(S0,0;u.t)=t$ for a constant update, and transport
of that equation along itself gives $t=t$.  Symmetry, transitivity and
term congruence then follow by transport.

Every induction below is an application of EI to one displayed
equation, with its other parameters fixed.  Finite-valued case analyses
are abbreviations for repeated NC and substitution of already proved
equations.  They do not invoke induction on a disjunction, a quantified
formula, or a simultaneous conjunction of equations.

\subsection{A finite orbit that exchanges zero and one}

Define the restricted selector and zero test
\[
 C^+(n,a,b):=R(n,a;u.b),\qquad Z(n):=C^+(n,0,0).
\]
Computation gives
\[
 C^+(0,a,b)=Sa,\quad C^+(Sn,a,b)=b,\qquad
 Z(0)=1,\quad Z(Sn)=0.
\]
The superscript emphasizes that this is not yet the ordinary selector.
Its zero branch adds one to its supplied value.

Put
\[
 \epsilon(n):=R(n,0;u.Z(u)),\qquad
 B(a,b):=Z\bigl(C^+(a,0,Z(b))\bigr).
\]
Thus $\epsilon(0)=1$ and $\epsilon(Sn)=Z(\epsilon(n))$.
Numerical cases on its counter and on the argument of $Z$ show that
$\epsilon(n)$ is always one of the two Boolean values.  On Boolean
arguments, $B$ has the conjunction table.  Consequently
\[
 \epsilon(SSn)=\epsilon(n).
 \label{ss:parity-period}
\]
No assertion about even or odd representations of arbitrary numbers is
needed for this equation.

Define
\[
 q(u):=B\bigl(\epsilon(u),Z(Z(u))\bigr),\qquad
 A(n):=R(n,1;u.q(u)),\qquad
 O(n):=Z\bigl(\epsilon(A(n))\bigr).
\]
Finite computation gives $q(2)=1$ and $q(1)=q(0)=0$.  Hence
\[
 A(0)=2,\quad A(1)=1,\quad A(SSn)=0;
 \qquad O(0)=0,\quad O(1)=1,\quad O(SSn)=0.
 \label{ss:one-indicator}
\]
For $A(SSn)=0$, the base is the two-step computation; its induction
step applies $q$ and uses $q(0)=0$.  The three equations for $O$ then
follow by substitution and finite computation.

Now set
\[
 j(u):=C^+\bigl(O(u),C^+(u,0,u),0\bigr),\qquad
 \tau(n):=R(n,0;u.j(u)).
\]
The preceding equations give
\[
 j(0)=2,\qquad j(1)=0,\qquad j(SSn)=SSSn.
\]
It follows that
\begin{equation}
 \tau(0)=1,\qquad \tau(1)=0,\qquad \tau(SSn)=SSn.
 \label{ss:swap}
\end{equation}
The third equation follows by EI: its base is
$\tau(2)=j(\tau(1))=j(0)=2$; its step uses the displayed equation
for $j(SSn)$.  Two numerical cases on an arbitrary input now prove
\begin{equation}
                         \tau(\tau(n))=n.
 \label{ss:swap-involution}
\end{equation}
This proof requires neither predecessor nor successor injectivity.

\subsection{Zero-start iteration by conjugation}

For every fixed scalar update $h$, define
\begin{equation}
 K(n;u.h):=
 \tau\!\left(R\left(n,0;
       v.\tau\bigl(h[\tau(v)/u]\bigr)\right)\right).
 \label{ss:zero-iterator}
\end{equation}
Bound variables are renamed before this finite substitution.  Its
computation equations are
\begin{equation}
 K(0;u.h)=0,\qquad
 K(Sn;u.h)=h[K(n;u.h)/u].
 \label{ss:k-computation}
\end{equation}
The base is $\tau(1)=0$.  At a successor, the two outer occurrences
of $\tau$ cancel by \eqref{ss:swap-involution}.
This gives a genuine zero-start scalar iterator for arbitrary fixed
updates, not merely for the finite updates used to construct $\tau$.

\subsection{Affine arithmetic before predecessor}

The shifted sum
\[
 T(x,y):=R(y,x;u.Su)
\]
satisfies $T(x,0)=Sx$ and $T(x,Sy)=S T(x,y)$.
Equation inductions give
\begin{equation}
 T(0,y)=Sy,\quad T(Sx,y)=S T(x,y),\quad
 T(x,y)=T(y,x),\quad
 T(T(x,y),z)=T(x,T(y,z)).
 \label{ss:shifted-semigroup}
\end{equation}
For example, the first two identities are proved by induction on $y$;
commutativity is then proved by induction on $y$, and associativity by
induction on $z$.  In particular, we have a commutative semigroup, not
an addition operation with zero as identity.

Define
\[
 d(n):=K(n;u.SSu),\qquad D(n):=T(n,n),\qquad
 \Theta(n):=T(T(n,n),n).
\]
Finite computation, \eqref{ss:shifted-semigroup}, and EI give
\begin{align}
 d(0)&=0,& d(Sn)&=SSd(n),&D(n)&=Sd(n),\notag\\
 D(0)&=1,&D(Sn)&=SSD(n),\notag\\
 \Theta(0)&=2,&\Theta(Sn)&=SSS\Theta(n),\notag\\
 \Theta(T(x,y))&=T(\Theta(x),\Theta(y)).
 \label{ss:affine-laws}
\end{align}
The last equality is simply reassociation and permutation of the six
arguments after expanding the three copies of $T(x,y)$.

Let $N(n):=Z(\epsilon(n))$.  The parity equations imply
\begin{equation}
 \epsilon(D(n))=0,\quad
 \epsilon(\Theta(n))=\epsilon(n),\quad
 N(\Theta(n))=N(n),\quad N(Sn)=Z(N(n)).
 \label{ss:affine-parity}
\end{equation}
For the first equation, induction follows the two successor increments
of $D$.  For the second, $\Theta$ advances by three successors, so both
sides toggle once.  The other equations follow by applying $Z$ and its
Boolean identities.

\subsection{A ramp and an affine factor map}

Define a three-state update and its zero-start orbit by
\[
 m(u):=C^+\bigl(u,0,C^+(\epsilon(u),1,0)\bigr),\qquad
 M(n):=K(n;u.m(u)).
\]
Only the finite values
\[
 m(0)=1,\qquad m(1)=2,\qquad m(2)=0
\]
will be used.  In particular, no general division algorithm or
remainder theorem is assumed.

Set
\[
 F(n):=K(n;u.T(u,M(u))).
\]
The central literal-equation induction is
\begin{equation}
                          M(F(n))=N(n).
 \label{ss:ramp-invariant}
\end{equation}
At zero, both sides are zero.  If $N(n)=0$, the next $F$ state is
$T(F(n),0)=SF(n)$, so its $M$ value is $m(0)=1=N(Sn)$.
If $N(n)=1$, the next state is $T(F(n),1)=SSF(n)$, whose $M$ value
is $m(m(1))=m(2)=0=N(Sn)$.  These are Boolean cases licensed by
the preceding definitions and NC.

It follows that
\begin{equation}
 F(Sn)=T(F(n),N(n)),\qquad F(SSn)=SSSF(n).
 \label{ss:ramp-step}
\end{equation}
The second equation uses the two Boolean cases of $N(n)$ and
$N(Sn)=Z(N(n))$.  A further induction, using the two-step increments
of $D$, gives
\begin{equation}
                          S F(D(n))=\Theta(n).
 \label{ss:ramp-double}
\end{equation}
Its base is $S F(1)=2=\Theta(0)$, and its successor step adds three
successors to both sides.

The second scaling identity is
\begin{equation}
 F(\Theta(n))=T(\Theta(F(n)),N(n)).
 \label{ss:ramp-triple}
\end{equation}
Here is a detailed induction check.  The zero case is $F(2)=3=T(2,0)$.
Write $X=\Theta(F(n))$ and $b=N(n)$.  After using the induction
equation, \eqref{ss:ramp-step}, and
$N(\Theta(n))=N(n)$, the new left side becomes
\[
                      S^3T(T(X,b),b).
\]
Using $F(Sn)=T(F(n),b)$ and the homomorphism law for $\Theta$, the
new right side becomes
\[
                    T(T(X,\Theta(b)),Z(b)).
\]
At $b=0$ both terms reduce to $S^5X$; at $b=1$ both reduce to $S^7X$,
using $\Theta(0)=2$ and $\Theta(1)=5$.  This is one equation induction
with a Boolean case split, not induction with a varying parameter.

Now define
\[
                   H(u):=C^+(\epsilon(u),F(u),u).
\]
The essential factor equations are
\begin{equation}
 H(0)=0,\qquad H(D(u))=\Theta(u),\qquad
 H(\Theta(u))=\Theta(H(u)).
 \label{ss:factor-map}
\end{equation}
The first is a computation.  The second follows from
$\epsilon(D(u))=0$ and \eqref{ss:ramp-double}.
For the third, $\epsilon(\Theta(u))=\epsilon(u)$.
If that Boolean value is $1$, both sides are $\Theta(u)$.
If it is $0$, then $N(u)=1$ and \eqref{ss:ramp-triple} gives
\[
 H(\Theta(u))=S F(\Theta(u))
   =S T(\Theta(F(u)),1)
   =S^3\Theta(F(u))
   =\Theta(SF(u))=\Theta(H(u)).
\]
The restricted selector is exactly sufficient: the transformed branch
is a successor, and the retained branch is the unmodified state.

\subsection{A strict comparator without general initialization}

Define
\[
 E(x):=K(x;u.D(u)),\qquad
 V(x,y):=\epsilon\bigl(R(y,d(E(x));u.H(u))\bigr).
\]
We prove the three equations
\begin{equation}
 V(x,0)=0,\qquad V(0,Sy)=1,\qquad
 V(Sx,Sy)=V(x,y).
 \label{ss:strict-comparator}
\end{equation}
The standard-model meaning of $V$ is the indicator of $x<y$, but
that interpretation is not used in the proof.

At zero, the recursor state is $Sd(E(x))=D(E(x))$, so the first
equation follows from \eqref{ss:affine-parity}.
For the second, $E(0)=d(0)=0$ and the first update of the initial
state $1=D(0)$ is $H(1)=\Theta(0)=2$.  Moreover $H(2)=2$.
Equation induction on $y$ therefore gives
$R(Sy,0;u.H(u))=2$, and $\epsilon(2)=1$.

For the third equation, prove by EI on $y$, with $x$ fixed,
\begin{equation}
 R(Sy,d(E(Sx));u.H(u))
     =\Theta\bigl(R(y,d(E(x));u.H(u))\bigr).
 \label{ss:trajectory-scaling}
\end{equation}
At $y=0$, the left side is
\[
 H(Sd(E(Sx)))=H(D(E(Sx)))
   =\Theta(E(Sx))=\Theta(D(E(x))),
\]
which is the right side's initial value.  The induction step applies
$H$ and then uses $H\Theta=\Theta H$.
Applying $\epsilon$ to \eqref{ss:trajectory-scaling} proves the third
equation of \eqref{ss:strict-comparator}.
This argument does not use an initializer-shifting rotation rule for
$R$; such a rule would require precisely the missing predecessor.

\subsection{Saturation constructs predecessor}

Put
\[
 B_m(u):=C^+(Z(V(u,m)),u,u),\qquad
 W(t,m):=K(t;u.B_m(u)).
\]
Because $V$ is Boolean, $B_m(u)$ is $Su$ when $V(u,m)=1$ and is
$u$ when $V(u,m)=0$.  Direct numerical cases give
\[
 C^+(e,Sa,Sb)=S C^+(e,a,b).
\]
Together with \eqref{ss:strict-comparator}, this yields
\[
 B_{Sm}(Su)=S B_m(u),\qquad B_{Sm}(0)=1.
\]
Equation induction on $t$, with $m$ fixed, now proves
\begin{equation}
                         W(St,Sm)=S W(t,m).
 \label{ss:saturated-shift}
\end{equation}
Its base is $B_{Sm}(0)=1$; its step applies the first of the two
preceding equations.  Two separate equation inductions give
\[
                       W(n,n)=n,\qquad V(n,n)=0.
\]
Consequently
\begin{equation}
                         W(Sn,n)=B_n(n)=n.
 \label{ss:extra-step}
\end{equation}

Define the finite scalar term
\begin{equation}
 P(m):=K\bigl(m;u.C^+(Z(V(Su,m)),u,u)\bigr).
 \label{ss:predecessor}
\end{equation}
The zero equation is immediate.  At parameter $m=Sn$, the update in
\eqref{ss:predecessor} agrees pointwise with $B_n$, since
$V(Su,Sn)=V(u,n)$.  EI on a separate iteration counter $t$, keeping
$n$ fixed, therefore proves equality of the two zero-start trajectories.
Only afterwards substitute $t=Sn$.  Equation \eqref{ss:extra-step}
then gives
\begin{equation}
                         P(0)=0,\qquad P(Sn)=n.
 \label{ss:predecessor-equations}
\end{equation}
This explicitly supplies the required body-extensionality argument;
updates with different syntax were not silently identified.

\subsection{Recovery of the ordinary scalar presentation}

For every fixed initialized update define
\begin{equation}
 I(n,a;u.h):=
 P\!\left(R\left(n,a;v.S\bigl(h[P(v)/u]\bigr)\right)\right).
 \label{ss:initialized-compiler}
\end{equation}
Using \eqref{ss:predecessor-equations} in the two computation cases
proves
\[
 I(0,a;u.h)=a,\qquad
 I(Sn,a;u.h)=h[I(n,a;u.h)/u].
\]
The ordinary selector is now also available, for example as
\[
 C(n,a,b):=P(C^+(n,a,Sb)).
\]
Its endpoints are exactly $a$ and $b$.  All ordinary logical gates and
selectors can therefore be reconstructed after the arithmetic
bootstrap.  There is no circular appeal to full intuitionistic logic
or an already available unrestricted selector in the predecessor proof.

Conversely, in the ordinary scalar presentation define
\[
                    R(n,a;u.h):=I(n,Sa;u.h).
\]
These descriptions satisfy \eqref{ss:r0}--\eqref{ss:rs}.  The two
numerical translations preserve zero and successor, and each translated
recursor agrees with the original by EI on its counter: their zero
values agree and their step functions agree.  With passive parameters
fixed, this gives the termwise identities required for the two
interpretation composites.  The domain is unchanged.  Combining this
identity-domain equivalence with the established full scalar-iPRA
bridge gives equiconsistency and bi-interpretability with iPRA.

\paragraph{Scope of the result.}
The reduction removes the independent successor constructor and the
counter variable from the successor-seeded recursor's update language.
It does not establish an absolute minimum for axiomatizations or a
checker-size improvement by itself.  A complete checker comparison
must still count every syntax, substitution, scoping and proof rule.
